\documentclass[11pt]{article}
\usepackage[a4paper,margin=1in]{geometry}
\usepackage{amsmath,amssymb,amsthm}
\usepackage[T1]{fontenc}
\usepackage{lmodern}
\usepackage[hidelinks]{hyperref}
\setlength{\emergencystretch}{3em}

\theoremstyle{plain}
\newtheorem{theorem}{Theorem}
\newtheorem{lemma}[theorem]{Lemma}
\theoremstyle{definition}
\newtheorem{definition}[theorem]{Definition}
\theoremstyle{remark}
\newtheorem{remark}[theorem]{Remark}

\newcommand{\SP}{\mathcal S}

\title{A Counterexample to Conjecture 00000008256:\\
a Maximal Condorcet Domain Above the Stated Bound}
\author{%
  Feiyu\thanks{Independent researcher.}
  \and
  Claude (Opus 5.5)\thanks{Anthropic.}%
}
\date{2026-10-03}

\begin{document}
\maketitle

\begin{abstract}
The first claim of Conjecture 00000008256 asserts that every maximal Condorcet
domain on $n$ alternatives has at most $2^{n-1}\,n!/2^{n(n-1)/2}$ linear orders.
For $n=3$ this bound equals $3$. The single-peaked domain on three alternatives,
consisting of the four orders in which the middle alternative is not ranked last, is
a maximal Condorcet domain with $4$ orders. The refutation is checked in Lean~4
against Mathlib.
\end{abstract}

\section{The conjecture as stated}

The original text of conjecture 00000008256 reads:

\begin{quote}
Definition: The maximal subdomains of the Condorcet domain for voting paradoxes: the
maximal preference domains cond\_max admitting Condorcet winners. Conjecture: The
size of cond\_max (a maximal Condorcet domain) is at most the Fishburn type explicit
bound $2^{n-1} (n!)/2^{n(n-1)/2}$ (a Fishburn bound formula); the tightness of the
bound (domains attaining it) is given by white-list domains (a white-list tight
domain law); the count of maximal domains (the number of maximal domains at $n = 4$)
is exactly $7$ (a four-alternative seven-domain law); and the permutation-symmetric
maximization structure of maximal domains is characterized by the lattice of
separable preferences (a separable lattice characterization law).
\end{quote}

The Chinese version in \texttt{conjectures/00000008256.md} writes the bound as
$2^{n-1}\cdot(n!)/2^{n(n-1)/2}$.

\section{Reading of the statement}

The conjecture is a conjunction of four claims. We refute the first one, so the
conjecture is false whatever the other three claims mean.

Fix $n$ alternatives. A \emph{preference} is a strict linear order on the
alternatives, and a \emph{domain} is a set $D$ of preferences. A \emph{profile} over
$D$ with $m$ voters assigns to each voter $i\in\{1,\dots,m\}$ a preference
$\succ_i\in D$; different voters may have the same preference.

\begin{definition}
Let $(\succ_i)_{i\le m}$ be a profile.
\begin{enumerate}
\item Alternative $a$ \emph{beats} $b$ if more than half of the voters prefer $a$ to
  $b$: $\#\{i: a\succ_i b\}>m/2$.
\item $a$ is a \emph{Condorcet winner} if it beats every other alternative.
\item A domain $D$ \emph{admits Condorcet winners} if every profile over $D$ with an
  odd number of voters has a Condorcet winner.
\item $D$ is a \emph{maximal Condorcet domain} if it admits Condorcet winners and
  every domain $D'\supseteq D$ that admits Condorcet winners equals $D$.
\end{enumerate}
\end{definition}

Restricting to an odd number of voters is the standard convention: it rules out
ties, which otherwise block a strict majority winner even in very small domains
(two voters with opposite preferences). The first claim then reads:
\begin{equation}\label{eq:bound}
|D|\ \le\ B(n):=\frac{2^{n-1}\,n!}{2^{n(n-1)/2}}
\quad\text{for every maximal Condorcet domain $D$ on $n$ alternatives.}
\tag{B}
\end{equation}
If ``the size of cond\_max'' is read instead as the largest size of a Condorcet
domain, the claim is refuted by the same example, since that largest size is at least
the size of any maximal Condorcet domain.

\section{Result}

For $n=3$ we have $B(3)=2^2\cdot 6/2^3=3$. Write the alternatives as $0,1,2$ and a
preference as the list of alternatives from best to worst, so that $102$ means
$1\succ0\succ2$. Let
\[
\SP=\{012,\ 102,\ 120,\ 210\}
\]
be the orders in which the middle alternative $1$ is not ranked last. These are
exactly the orders that are single-peaked with respect to the axis $0<1<2$.

\begin{theorem}\label{thm:main}
$\SP$ is a maximal Condorcet domain on three alternatives with $|\SP|=4>3=B(3)$.
Hence (B) fails, and Conjecture 00000008256 is false.
\end{theorem}

\section{Proof}

\begin{lemma}\label{lem:black}
$\SP$ admits Condorcet winners.
\end{lemma}

\begin{proof}
This is Black's median voter theorem for three alternatives; we give the direct
argument. Take a profile over $\SP$ with $m$ voters, $m$ odd. Let $N_0$ and $N_2$
be the numbers of voters ranking $0$, respectively $2$, first.

If $N_0>m/2$, then $0$ is a Condorcet winner, since every voter ranking $0$ first
prefers $0$ to each other alternative. Likewise, if $N_2>m/2$, then $2$ is a
Condorcet winner.

Otherwise $N_0\le m/2$ and $N_2\le m/2$, hence $N_0<m/2$ and $N_2<m/2$ since $m$ is
odd. The only order in $\SP$ with $0$ on top is $012$, and in the other three orders
$102,120,210$ the alternative $1$ is above $0$. So exactly $m-N_0>m/2$ voters prefer
$1$ to $0$. Symmetrically, the only order in $\SP$ with $2$ on top is $210$, and in
$012,102,120$ the alternative $1$ is above $2$. So $m-N_2>m/2$ voters prefer $1$ to
$2$. Hence $1$ is a Condorcet winner.
\end{proof}

\begin{lemma}\label{lem:max}
$\SP$ is a maximal Condorcet domain.
\end{lemma}

\begin{proof}
By Lemma~\ref{lem:black}, $\SP$ admits Condorcet winners. Let $D'\supseteq\SP$
admit Condorcet winners, and suppose $D'$ contains an order outside $\SP$. The only
orders outside $\SP$ are $021$ and $201$.
\begin{itemize}
\item If $021\in D'$, consider the profile $(021,\,210,\,102)$ over $D'$. Here $0$
  beats $2$ (voters 1 and 3), $2$ beats $1$ (voters 1 and 2), and $1$ beats $0$
  (voters 2 and 3).
\item If $201\in D'$, consider the profile $(201,\,012,\,120)$ over $D'$. Here $0$
  beats $1$ (voters 1 and 2), $1$ beats $2$ (voters 2 and 3), and $2$ beats $0$
  (voters 1 and 3).
\end{itemize}
In either case the majority relation is a cycle, so every alternative is beaten by
some other one and there is no Condorcet winner, a contradiction. Hence $D'=\SP$.
\end{proof}

\begin{proof}[Proof of Theorem~\ref{thm:main}]
By Lemma~\ref{lem:max}, $\SP$ is a maximal Condorcet domain, and $|\SP|=4>3=B(3)$.
\end{proof}

\section{Formalization}

The Lean~4 project in \texttt{lean/} (file \texttt{lean/Submission/Basic.lean})
formalizes the definitions above and proves Theorem~\ref{thm:main} against Mathlib.
\begin{itemize}
  \item A preference on \texttt{Fin n} is encoded by its ranking
    \texttt{r : Equiv.Perm (Fin n)}, where \texttt{r a} is the position of $a$
    (with $0$ the top) and $a$ is preferred to $b$ iff \texttt{r a < r b}. This
    encodes every strict linear order exactly once.
  \item \texttt{Beats v a b} is \texttt{m < 2 * \#\{i | v i a < v i b\}} for a
    profile \texttt{v : Fin m -> Ranking n}; \texttt{IsCondorcetWinner},
    \texttt{AdmitsCondorcetWinners} (all profiles with \texttt{Odd m}) and
    \texttt{IsMaximalCondorcetDomain} follow the definition above.
  \item \texttt{SizeBound} is (B), with the bound computed in $\mathbb R$.
  \item \texttt{singlePeaked} is $\SP$, defined as the rankings with
    \texttt{r 1} $\ne 2$; \texttt{card\_singlePeaked} gives $|\SP|=4$.
  \item \texttt{admitsCondorcetWinners\_singlePeaked} is Lemma~\ref{lem:black},
    proved for all odd $m$ by the counting argument above.
    \texttt{cycle\_of\_not\_mem} provides the cyclic three-voter profiles, by
    \texttt{decide}, and \texttt{isMaximalCondorcetDomain\_singlePeaked} is
    Lemma~\ref{lem:max}.
  \item \texttt{conjecture\_00000008256\_false} is $\lnot\,$\texttt{SizeBound}, and
    \texttt{conjecture\_00000008256\_false'} states
    $\lnot(\texttt{SizeBound}\wedge P\wedge Q\wedge R)$ for arbitrary propositions
    $P,Q,R$ standing for the other three claims.
\end{itemize}
The toolchain is \texttt{leanprover/lean4:v4.33.1}, Mathlib is pinned to
\texttt{v4.33.1}, and \texttt{lake build} succeeds. The project contains no
\texttt{sorry}, no \texttt{admit} and no \texttt{native\_decide}, and introduces no
axioms: \texttt{\#print axioms conjecture\_00000008256\_false} reports exactly
\texttt{propext}, \texttt{Classical.choice} and \texttt{Quot.sound}.

\section{Remarks}

\begin{remark}
For three alternatives and an odd number of voters, the majority relation is a
tournament, and a tournament on three vertices has a Condorcet winner exactly when
it is acyclic. So on three alternatives our notion coincides with the other common
definition of a Condorcet domain (the majority relation is transitive for every
profile with an odd number of voters), and $\SP$ is maximal in that sense too.
\end{remark}

\begin{remark}
The bound fails for every $n\ge3$, although only $n=3$ is formalized. The
single-peaked domain for an axis on $n$ alternatives has $2^{n-1}$ orders and admits
Condorcet winners by Black's theorem~\cite{black}, and it lies in some maximal
Condorcet domain. On the other hand $B(n)<2^{n-1}$, because
$n!<2^{n(n-1)/2}$ for $n\ge3$. Moreover $B(n+1)/B(n)=(n+1)/2^{n-1}$, so
$B(3)=B(4)=3$, $B(5)=15/8$, $B(6)=45/64$, and $B(n)<1$ for all $n\ge6$. Every
maximal Condorcet domain is nonempty, because a domain with a single order admits
Condorcet winners. So for $n\ge6$ every maximal Condorcet domain violates (B).
\end{remark}

\begin{remark}
We do not address the other three claims of the conjecture.
\end{remark}

\begin{thebibliography}{9}
\bibitem{black} D.~Black, \emph{On the rationale of group decision-making},
  J. Political Economy 56 (1948), 23--34.
\end{thebibliography}

\end{document}
